\SourcePage{205}{208}
\chapter{The semiclassical case}
\section{The wave function in the semiclassical case}

When the de Broglie wavelengths of the particles are small in comparison with
the characteristic dimensions $L$ that determine the conditions of a given
problem, the system has nearly classical properties. (This is analogous to
the passage from wave optics to geometrical optics as the wavelength tends to
zero.)

We shall now examine the properties of semiclassical systems in greater
detail. In the Schrödinger equation
\[
  \sum_a\frac{\hbar^2}{2m_a}\Delta_a\psi+(E-U)\psi=0
\]
make the formal substitution
\begin{equation}
  \psi=\exp\left(\frac i\hbar\sigma\right).
  \tag{46.1}
\end{equation}
The resulting equation for $\sigma$ is
\begin{equation}
  \sum_a\frac1{2m_a}(\GradAt{a}{\sigma})^2
  -\sum_a\frac{i\hbar}{2m_a}\Delta_a\sigma=E-U.
  \tag{46.2}
\end{equation}
In keeping with the assumption that the properties of the system are nearly
classical, seek $\sigma$ as the series
\begin{equation}
  \sigma=\sigma_0+\frac\hbar i\sigma_1+
  \left(\frac\hbar i\right)^2\sigma_2+\cdots,
  \tag{46.3}
\end{equation}
arranged in powers of $\hbar$.

Begin with the simplest case, one-dimensional motion of a single particle.
Equation (46.2) then reduces to
\begin{equation}
  \frac1{2m}\sigma'^2-\frac{i\hbar}{2m}\sigma''=E-U(x)
  \tag{46.4}
\end{equation}
(a prime denotes differentiation with respect to the coordinate $x$).

In the first approximation put $\sigma=\sigma_0$ and omit the term containing
$\hbar$ from the equation:
\[
  \frac1{2m}\sigma_0'^2=E-U(x).
\]
\SourcePage{206}{209}
It follows that
\[
  \sigma_0=\pm\int\sqrt{2m[E-U(x)]}\,dx.
\]
The integrand is precisely the classical momentum $p(x)$ of the particle,
expressed as a function of the coordinate. Defining $p(x)$ with the positive
sign before the square root, we have
\begin{equation}
  \sigma_0=\pm\int p\,dx,
  \qquad p=\sqrt{2m(E-U)},
  \tag{46.5}
\end{equation}
as expected from the limiting expression (6.1) for the wave
function.\footnote{As is well known, $\int p\,dx$ is the time-independent part
of the action. The complete mechanical action of the particle is
$S=-Et\pm\int p\,dx$. The term $Et$ is absent from $\sigma$ because we are
considering the time-independent wave function $\psi$.}

The neglect made in (46.4) is justified only when the second term on the
left-hand side is small compared with the first. Thus
$\hbar|\sigma''/\sigma'^2|\ll1$, or
\[
  \left|\frac d{dx}\left(\frac\hbar{\sigma'}\right)\right|\ll1.
\]
To first order, (46.5) gives $\sigma'=p$, so this condition can be written
\begin{equation}
  \left|\frac{d\lambdabar}{dx}\right|\ll1,
  \tag{46.6}
\end{equation}
where $\lambdabar=\lambda/2\pi$, while
$\lambda(x)=2\pi\hbar/p(x)$ is the particle's de Broglie wavelength expressed
as a function of $x$ through the classical function $p(x)$. We have thereby
obtained a quantitative semiclassical condition: over a distance comparable
with the wavelength itself, that wavelength must vary only slightly. The
approximation fails in regions of space where this condition is not met.

Condition (46.6) may also be put in another form. Since
\[
  \frac{dp}{dx}=\frac d{dx}\sqrt{2m(E-U)}=-\frac mp\frac{dU}{dx}=\frac{mF}{p},
\]
where $F=-dU/dx$ is the classical force acting on the particle in the external
field, introduction of this force gives
\begin{equation}
  \frac{m\hbar|F|}{p^3}\ll1.
  \tag{46.7}
\end{equation}
\SourcePage{207}{210}
This shows that the semiclassical approximation fails when the particle's
momentum is too small. In particular, it certainly fails near turning points,
namely points at which classical mechanics would have the particle stop and
then begin moving in the opposite direction. Such points are determined by
$p(x)=0$, or $E=U(x)$. As $p\to0$, the de Broglie wavelength tends to infinity
and plainly cannot be regarded as small.

It should be stressed, however, that condition (46.6) or (46.7) may by itself
be insufficient to justify the semiclassical approximation. It was obtained
by estimating different terms in the differential equation (46.4), and the
discarded term contains the highest derivative. What must actually be small
are successive terms in the expansion of the solution of this equation; the
smallness of the discarded term in the equation need not ensure this. For
example, suppose the solution for $\sigma(x)$ contains a term whose growth
with $x$ is nearly linear. The smallness of its second derivative does not
prevent that term from accumulating a large value over a sufficiently great
distance. Such a situation generally arises when the field extends over
distances large compared with the characteristic length $L$ on which it
changes appreciably (see the remark below concerning (46.11)). The
semiclassical approximation then cannot follow the wave function over large
distances.

We proceed to calculate the next term of (46.3). Terms of first order in
$\hbar$ in (46.4) give
\[
  \sigma_0'\sigma_1'+\frac12\sigma_0''=0,
  \qquad
  \sigma_1'=-\frac{\sigma_0''}{2\sigma_0'}=-\frac{p'}{2p}.
\]
Integration yields
\begin{equation}
  \sigma_1=-\frac12\ln p
  \tag{46.8}
\end{equation}
(the integration constant is omitted).

Substitution of this expression in (46.1) and (46.3) gives the wave function
\begin{equation}
  \psi=\frac C{\sqrt p}\exp\left(\frac i\hbar\int p\,dx\right)
  +\frac{C'}{\sqrt p}\exp\left(-\frac i\hbar\int p\,dx\right).
  \tag{46.9}
\end{equation}
The factor $1/\sqrt p$ has a simple interpretation. The probability of finding
the particle between $x$ and $x+dx$ is determined by $|\psi|^2$, and hence is
\SourcePage{208}{211}
principally proportional to $1/p$. This is just what one expects for a
``semiclassical particle,'' since in classical motion the time spent by a
particle in an interval $dx$ is inversely proportional to its velocity (or
momentum).

In classically inaccessible regions, where $E<U(x)$, the function $p(x)$ is
purely imaginary and the exponential arguments are real. The general form of
the wave-equation solution in such regions is
\begin{equation}
  \psi=\frac C{\sqrt{|p|}}\exp\left(-\frac1\hbar\int|p|\,dx\right)
  +\frac{C'}{\sqrt{|p|}}\exp\left(\frac1\hbar\int|p|\,dx\right).
  \tag{46.10}
\end{equation}
It must nevertheless be remembered that the accuracy of the semiclassical
approximation does not permit exponentially small terms to be retained in a
wave function against a background of exponentially large terms. In this
sense, retaining both terms in (46.10) simultaneously is generally
inadmissible.

Although higher-order terms in the wave function are usually unnecessary, we
shall also obtain the next term in (46.3), in order to point out certain
features concerning the accuracy of the semiclassical approximation.

The terms of order $\hbar^2$ in (46.4) give
\[
  \sigma_0'\sigma_2'+\frac12\sigma_1'^2+\frac12\sigma_1''=0,
\]
whence, on substituting (46.5) and (46.8) for $\sigma_0$ and $\sigma_1$,
\[
  \sigma_2'=\frac{p''}{4p^2}-\frac{3p'^2}{8p^3}.
\]
On integrating (the first term being integrated by parts) and introducing the
force $F=pp'/m$, we find
\[
  \sigma_2=\frac{mF}{4p^3}+\frac{m^2}{8}\int\frac{F^2}{p^5}\,dx.
\]
In the approximation under consideration the wave function is
\[
  \psi=\exp[(i/\hbar)\sigma]
  =\exp[(i/\hbar)\sigma_0+\sigma_1](1-i\hbar\sigma_2),
\]
or
\begin{equation}
  \psi=\frac{\operatorname{const}}{\sqrt p}
  \left[1-\frac{im\hbar F}{4p^3}
  -\frac{im^2\hbar}{8}\int\frac{F^2}{p^5}\,dx\right]
  \exp\left(\frac i\hbar\int p\,dx\right).
  \tag{46.11}
\end{equation}
The imaginary correction terms in the prefactor are equivalent to a real
correction in the phase of the wave function (that is, an addition to the
\SourcePage{209}{212}
integral $\hbar^{-1}\int p\,dx$ in its exponent). This correction is
proportional to $\hbar$, and is therefore of order $\lambda/L$.

The second and third terms in the square brackets in (46.11) must be small
compared with unity. For the former this requirement is identical to (46.7),
whereas for the latter an estimate of the integral leads to (46.7) only if
$F^2$ tends to zero sufficiently rapidly over distances of order $L$.
